\documentclass[UTF8,fontset=fandol,11pt,a4paper]{ctexart}

\usepackage{amsmath,amssymb}
\usepackage{booktabs}
\usepackage{array}
\usepackage{geometry}
\usepackage{graphicx}
\usepackage{xcolor}
\usepackage[numbers]{natbib}
\usepackage[hidelinks]{hyperref}

\geometry{left=2.6cm,right=2.6cm,top=2.5cm,bottom=2.5cm}
\setlength{\parskip}{0.35em}
\setlength{\parindent}{2em}

% 结果占位：正文中任何未完成的结论句都用 \TOPFILL 标记，避免手写数字。
\newcommand{\TOPFILL}[1]{\textcolor{red}{[待填：{#1}]}}

% 记忆消融宏（由 compare_conditions.py 生成；缺失时 \RESULT 展开为 --）
% 占位：跑批完成后由 compare_conditions.py 覆盖


% 所有实验结果单元格只允许通过 \RESULT{臂}{指标} 取值；宏由
% tools/bench/analysis/render_tables.py 从 artifacts/bench 生成。
% 产物缺失时展开为 --，绝不回落到 0。
\newcommand{\RESULT}[2]{%
  \ifcsname res@#1@#2\endcsname
    \csname res@#1@#2\endcsname
  \else
    --%
  \fi}
\IfFileExists{tables/macros.tex}{% 由 tools/bench/analysis/render_tables.py 生成，请勿手改。
% 数据源：shard0.json+shard1.json+shard2.json+shard3.json（分片，部分产物）
% 生成时间（UTC）：2026-11-27T16:17:25+00:00

\expandafter\def\csname res@collaborative@attempted\endcsname{4}
\expandafter\def\csname res@collaborative@balacc\endcsname{0.500}
\expandafter\def\csname res@collaborative@base_rate\endcsname{0.500}
\expandafter\def\csname res@collaborative@coverage\endcsname{0.033}
\expandafter\def\csname res@collaborative@degraded_rate\endcsname{0.000}
\expandafter\def\csname res@collaborative@degraded_roles\endcsname{0}
\expandafter\def\csname res@collaborative@failures\endcsname{0}
\expandafter\def\csname res@collaborative@fn\endcsname{2}
\expandafter\def\csname res@collaborative@fp\endcsname{0}
\expandafter\def\csname res@collaborative@latency_s\endcsname{98.2}
\expandafter\def\csname res@collaborative@leakage\endcsname{0}
\expandafter\def\csname res@collaborative@lift\endcsname{--}
\expandafter\def\csname res@collaborative@n\endcsname{4}
\expandafter\def\csname res@collaborative@null_risk\endcsname{0}
\expandafter\def\csname res@collaborative@partial_rate\endcsname{0.000}
\expandafter\def\csname res@collaborative@planned\endcsname{120}
\expandafter\def\csname res@collaborative@precision\endcsname{--}
\expandafter\def\csname res@collaborative@precision_ci\endcsname{--}
\expandafter\def\csname res@collaborative@recall\endcsname{0.000}
\expandafter\def\csname res@collaborative@recall_ci\endcsname{[0.000, 0.658]}
\expandafter\def\csname res@collaborative@tn\endcsname{2}
\expandafter\def\csname res@collaborative@tokens\endcsname{24597}
\expandafter\def\csname res@collaborative@tp\endcsname{0}
\expandafter\def\csname res@collaborative@trigger\endcsname{0.000}
\expandafter\def\csname res@collaborative@unsupported\endcsname{0}
\expandafter\def\csname res@collaborative@valid\endcsname{4}
\expandafter\def\csname res@collaborative@verify_failed\endcsname{0}
\expandafter\def\csname res@collaborative@verify_pass\endcsname{1.000}
\expandafter\def\csname res@independent@attempted\endcsname{4}
\expandafter\def\csname res@independent@balacc\endcsname{0.500}
\expandafter\def\csname res@independent@base_rate\endcsname{0.500}
\expandafter\def\csname res@independent@coverage\endcsname{0.033}
\expandafter\def\csname res@independent@degraded_rate\endcsname{0.000}
\expandafter\def\csname res@independent@degraded_roles\endcsname{0}
\expandafter\def\csname res@independent@failures\endcsname{0}
\expandafter\def\csname res@independent@fn\endcsname{2}
\expandafter\def\csname res@independent@fp\endcsname{0}
\expandafter\def\csname res@independent@latency_s\endcsname{82.8}
\expandafter\def\csname res@independent@leakage\endcsname{0}
\expandafter\def\csname res@independent@lift\endcsname{--}
\expandafter\def\csname res@independent@n\endcsname{4}
\expandafter\def\csname res@independent@null_risk\endcsname{0}
\expandafter\def\csname res@independent@partial_rate\endcsname{0.000}
\expandafter\def\csname res@independent@planned\endcsname{120}
\expandafter\def\csname res@independent@precision\endcsname{--}
\expandafter\def\csname res@independent@precision_ci\endcsname{--}
\expandafter\def\csname res@independent@recall\endcsname{0.000}
\expandafter\def\csname res@independent@recall_ci\endcsname{[0.000, 0.658]}
\expandafter\def\csname res@independent@tn\endcsname{2}
\expandafter\def\csname res@independent@tokens\endcsname{20362}
\expandafter\def\csname res@independent@tp\endcsname{0}
\expandafter\def\csname res@independent@trigger\endcsname{0.000}
\expandafter\def\csname res@independent@unsupported\endcsname{0}
\expandafter\def\csname res@independent@valid\endcsname{4}
\expandafter\def\csname res@independent@verify_failed\endcsname{0}
\expandafter\def\csname res@independent@verify_pass\endcsname{1.000}
\expandafter\def\csname res@label@base_strict\endcsname{0.341}
\expandafter\def\csname res@label@base_y20\endcsname{0.341}
\expandafter\def\csname res@label@n\endcsname{2646}
\expandafter\def\csname res@label@n_agree\endcsname{2646}
\expandafter\def\csname res@label@n_hs0\endcsname{2341}
\expandafter\def\csname res@label@n_hs1\endcsname{305}
\expandafter\def\csname res@label@n_strict\endcsname{2646}
\expandafter\def\csname res@label@p_strict_hs0\endcsname{0.265}
\expandafter\def\csname res@label@p_strict_hs1\endcsname{0.921}
\expandafter\def\csname res@label@p_y_hs0\endcsname{0.265}
\expandafter\def\csname res@label@p_y_hs1\endcsname{0.921}
\expandafter\def\csname res@mcnemar_independent@n\endcsname{4}
\expandafter\def\csname res@mcnemar_independent@only_base\endcsname{0}
\expandafter\def\csname res@mcnemar_independent@only_collab\endcsname{0}
\expandafter\def\csname res@mcnemar_independent@p\endcsname{--}
\expandafter\def\csname res@mcnemar_independent@p_holm\endcsname{--}
\expandafter\def\csname res@mcnemar_independent@verdict\endcsname{无不一致对}
\expandafter\def\csname res@mcnemar_single@n\endcsname{4}
\expandafter\def\csname res@mcnemar_single@only_base\endcsname{0}
\expandafter\def\csname res@mcnemar_single@only_collab\endcsname{0}
\expandafter\def\csname res@mcnemar_single@p\endcsname{--}
\expandafter\def\csname res@mcnemar_single@p_holm\endcsname{--}
\expandafter\def\csname res@mcnemar_single@verdict\endcsname{无不一致对}
\expandafter\def\csname res@meta@base_url\endcsname{http://100.121.213.4:3065/v1}
\expandafter\def\csname res@meta@generated_at\endcsname{2026-11-27T16:16:11.493254Z}
\expandafter\def\csname res@meta@model\endcsname{deepseek-v4.1-flash}
\expandafter\def\csname res@meta@n_planned\endcsname{120}
\expandafter\def\csname res@meta@n_records\endcsname{4}
\expandafter\def\csname res@meta@n_scored\endcsname{4}
\expandafter\def\csname res@meta@n_shards\endcsname{4}
\expandafter\def\csname res@meta@partial\endcsname{是}
\expandafter\def\csname res@meta@seed\endcsname{11}
\expandafter\def\csname res@meta@source\endcsname{shard0.json+shard1.json+shard2.json+shard3.json（分片，部分产物）}
\expandafter\def\csname res@meta@threshold\endcsname{high}
\expandafter\def\csname res@placebo_collaborative@draws\endcsname{0}
\expandafter\def\csname res@placebo_collaborative@null_mean\endcsname{--}
\expandafter\def\csname res@placebo_collaborative@null_q95\endcsname{--}
\expandafter\def\csname res@placebo_collaborative@obs_lift\endcsname{--}
\expandafter\def\csname res@placebo_collaborative@p\endcsname{--}
\expandafter\def\csname res@placebo_collaborative@verdict\endcsname{未执行（无触发或无真值）}
\expandafter\def\csname res@placebo_independent@draws\endcsname{0}
\expandafter\def\csname res@placebo_independent@null_mean\endcsname{--}
\expandafter\def\csname res@placebo_independent@null_q95\endcsname{--}
\expandafter\def\csname res@placebo_independent@obs_lift\endcsname{--}
\expandafter\def\csname res@placebo_independent@p\endcsname{--}
\expandafter\def\csname res@placebo_independent@verdict\endcsname{未执行（无触发或无真值）}
\expandafter\def\csname res@placebo_single@draws\endcsname{0}
\expandafter\def\csname res@placebo_single@null_mean\endcsname{--}
\expandafter\def\csname res@placebo_single@null_q95\endcsname{--}
\expandafter\def\csname res@placebo_single@obs_lift\endcsname{--}
\expandafter\def\csname res@placebo_single@p\endcsname{--}
\expandafter\def\csname res@placebo_single@verdict\endcsname{未执行（无触发或无真值）}
\expandafter\def\csname res@single@attempted\endcsname{4}
\expandafter\def\csname res@single@balacc\endcsname{0.500}
\expandafter\def\csname res@single@base_rate\endcsname{0.500}
\expandafter\def\csname res@single@coverage\endcsname{0.033}
\expandafter\def\csname res@single@degraded_rate\endcsname{0.000}
\expandafter\def\csname res@single@degraded_roles\endcsname{0}
\expandafter\def\csname res@single@failures\endcsname{0}
\expandafter\def\csname res@single@fn\endcsname{2}
\expandafter\def\csname res@single@fp\endcsname{0}
\expandafter\def\csname res@single@latency_s\endcsname{73.6}
\expandafter\def\csname res@single@leakage\endcsname{0}
\expandafter\def\csname res@single@lift\endcsname{--}
\expandafter\def\csname res@single@n\endcsname{4}
\expandafter\def\csname res@single@null_risk\endcsname{0}
\expandafter\def\csname res@single@partial_rate\endcsname{0.000}
\expandafter\def\csname res@single@planned\endcsname{120}
\expandafter\def\csname res@single@precision\endcsname{--}
\expandafter\def\csname res@single@precision_ci\endcsname{--}
\expandafter\def\csname res@single@recall\endcsname{0.000}
\expandafter\def\csname res@single@recall_ci\endcsname{[0.000, 0.658]}
\expandafter\def\csname res@single@tn\endcsname{2}
\expandafter\def\csname res@single@tokens\endcsname{18954}
\expandafter\def\csname res@single@tp\endcsname{0}
\expandafter\def\csname res@single@trigger\endcsname{0.000}
\expandafter\def\csname res@single@unsupported\endcsname{0}
\expandafter\def\csname res@single@valid\endcsname{4}
\expandafter\def\csname res@single@verify_failed\endcsname{0}
\expandafter\def\csname res@single@verify_pass\endcsname{1.000}
}{}

\title{\bfseries 多主体协同带来判别增量吗？\\[3pt]
\large 面向跨境金融风险的可审计预警系统 CBR-MAW 的消融研究}
\author{CBR-MAW 研究组}
\date{2026 年 11 月}

\begin{document}
\sloppy
\emergencystretch=3em
\maketitle

\begin{abstract}
跨境金融风险预警既要求判别力，也要求可审计、可回放且无前视偏差。多智能体协同
被广泛假设为能够改善复杂判断，但在“同一数据、同一模型、同一工具”的严格对照下，
它是否真的改变风险定级与预警判别力，仍缺少可复现的证据。本文提出并实现
CBR-MAW：一个纯客户端、三进程的跨境金融风险多主体预警系统。Dart 桌面外壳承载
交互与编排，Go \texttt{collector} 负责多源采集与清洗，Rust \texttt{kernel} 作为
唯一写入者承载 PIT 时序、三库记忆与确定性数值工具；内环 R0--R5 提供证据共享、
交叉验证、分歧分类、仲裁与六类核验，外环以 MI/LLR/HLR 三库记忆和浅层/深层反思
形成跨时点闭环；四层 PIT 过滤与内容寻址的证据绑定约束全部数值与引文。我们给出
一套预注册的配对消融协议：在同一模型端点下比较 \texttt{collaborative}、
\texttt{single}、\texttt{independent} 三臂，报告 precision、recall、lift、
Wilson 置信区间、精确 McNemar、触发率、延迟与 token，并强制报告覆盖率、遮蔽计数、
标签置换安慰剂与标签循环性诊断。主要结果：在 \RESULT{meta}{n_scored} 个可评分配对时点上
（计划 \RESULT{meta}{n_planned}），协同臂均衡准确率相对单智能体臂差
\RESULT{concl}{balacc_delta_single}、相对独立臂差 \RESULT{concl}{balacc_delta_indep}；
经 Holm 校正的配对检验判定为「\RESULT{concl}{vs_single}」与「\RESULT{concl}{vs_indep}」
（$n\ge 30$ 是否满足：\RESULT{concl}{enough_n}）。预测向量一致性：
\RESULT{concl}{identical}。成本上协同臂 token 为单智能体的 \RESULT{concl}{token_ratio} 倍、
耗时为 \RESULT{concl}{latency_ratio} 倍。本文同样报告负结果与失效条件，并披露
最小覆盖率 \RESULT{concl}{min_coverage}、最高降级率 \RESULT{concl}{max_degraded}、
是否存在部分结论标记（\RESULT{concl}{partial_any}）。
所有实验数字均由脚本从真实产物生成，正文不手工填写结果。

\vspace{0.6em}
\noindent\textbf{关键词：}跨境金融风险；多智能体系统；可审计人工智能；
时点一致性；消融研究；金融压力指数
\end{abstract}

\vspace{1.2em}
\begin{center}
{\large\bfseries Abstract}
\end{center}
\begin{quote}
Cross-border financial risk early warning requires discriminative power together with
auditability, replayability, and the absence of look-ahead bias. Multi-agent collaboration
is often assumed to improve complex judgments, yet it is unclear whether it changes risk
grading or warning discrimination under a strict control of identical data, identical model,
and identical tools. We present CBR-MAW, a purely client-side, three-process early-warning
system for cross-border financial risk. A Dart desktop shell carries interaction and
orchestration, a Go \texttt{collector} performs multi-source ingestion and cleaning, and a
Rust \texttt{kernel} acts as the single writer for point-in-time (PIT) series, a three-store
memory, and deterministic numerical tools. The inner R0--R5 loop provides evidence sharing,
cross-validation, conflict classification, arbitration, and six-class verification; the outer
loop closes the cross-time cycle with MI/LLR/HLR memories and shallow/deep reflection.
Four layers of PIT filtering and content-addressed evidence binding constrain every number and
citation. We report a pre-registered paired ablation that compares \texttt{collaborative},
\texttt{single}, and \texttt{independent} arms on the same model endpoint, with precision,
recall, lift, Wilson intervals, exact McNemar tests, trigger rate, latency, token cost,
coverage, masking, label-permutation placebo, and label-circularity diagnostics.
Results: over \RESULT{meta}{n_scored} scored paired time points (\RESULT{meta}{n_planned}
planned), the collaborative arm's balanced accuracy differs from the single-agent arm by
\RESULT{concl}{balacc_delta_single} and from the independent arm by
\RESULT{concl}{balacc_delta_indep}; Holm-corrected paired tests are
``\RESULT{concl}{vs_single}'' and ``\RESULT{concl}{vs_indep}'' respectively
($n\ge 30$: \RESULT{concl}{enough_n}). Prediction vectors: \RESULT{concl}{identical}.
Cost-wise the collaborative arm consumes \RESULT{concl}{token_ratio}$\times$ the tokens and
\RESULT{concl}{latency_ratio}$\times$ the wall time of the single-agent arm. Negative results
are reported alongside positive ones; minimum coverage \RESULT{concl}{min_coverage},
maximum degraded rate \RESULT{concl}{max_degraded}, partial-flag present
(\RESULT{concl}{partial_any}). Every experimental number in this paper is generated by scripts
from real artifacts; no result is typed by hand.
\end{quote}

\noindent\textbf{Keywords:} cross-border financial risk; multi-agent systems; auditable AI;
point-in-time consistency; ablation study; financial stress index

\section{引言}

\subsection{问题背景}

跨境金融风险具有三个同时存在的约束。第一，风险状态由多市场、多币种、多机构的
异构序列共同决定，单一指标难以覆盖信用、汇率、权益、债券、跨境流动与全球风险
等维度\citep{illing2006measuring,hollo2012ciss,cardarelli2011financial}。第二，监管与
机构场景要求结论可追溯：每个数值应当能回指到确定的输入、工具版本与时点，而不是
一段无法复核的自然语言\citep{schneier1999audit,parno2013pinocchio}。第三，历史回放
必须严格时点一致（point-in-time, PIT）：在评估日 $t$ 只能使用 $t$ 当日可得的信息，
否则回测会把未来数据泄漏进历史判断\citep{kaufman2012leakage,bailey2016overfitting}。

大语言模型（LLM）与多智能体系统为上述问题提供了新的工程路径。多智能体辩论、
角色分工与工具调用在事实性、推理与任务完成度上展现出潜力
\citep{du2023debate,liang2023mad,wu2023autogen,hong2023metagpt}；检索增强、长期记忆与
反思机制则被用于把跨时点经验带入当前判断
\citep{lewis2020rag,park2023generative,shinn2023reflexion,packer2023memgpt}。
然而，这些工作大多在问答、推理或开放式任务上评估，较少把“协同是否改变一个
带时点约束的金融风险定级”作为可证伪的对照问题；而金融预警恰好同时具备小样本、
强时间依赖、标签与指标同源等让评估极易自欺的特征。

\subsection{本文的问题与预注册假设}

本文把问题限定为一个可操作的对照问题：在相同数据快照、相同模型端点、相同工具
注册表与相同记忆配置下，把三主体从“独立分析”改为“完整协同”，是否改变风险等级
与预警判别力？需要付出多少延迟与 token 成本？过程指标（覆盖率、核验、遮蔽）如何
变化？

我们在实验前登记以下假设，正负结果都报告：
\begin{itemize}
  \item \textbf{H1（判别增量）}：协同臂的均衡准确率高于单智能体臂与独立臂；
        用双侧精确 McNemar 检验，
        并对两项比较做 Holm 校正，显著性水平 $\alpha=0.05$。
  \item \textbf{H2（成本）}：协同臂的平均端到端延迟与平均 token 消耗高于
        单智能体臂与独立臂。
  \item \textbf{H3（过程覆盖）}：协同臂的有效定级覆盖率与核验通过率不低于
        单臂；该假设只作描述性报告，不作显著性主张。
  \item \textbf{H0}：三臂在 H1 上无差异。
\end{itemize}
预注册并不保证结果符合预期；若 H1 被证伪，本文的价值在于把“协同无增量”的
条件、成本与失效机制记录清楚，而不是把负结果重新叙述成正结果。

\subsection{贡献}

本文的贡献限于“提出、实现与验证流程”，不主张尚未验证的收益：
\begin{enumerate}
  \item \textbf{架构}：提出并实现一个纯客户端、三进程的多主体跨境金融风险预警
        系统。系统把交互与编排、采集清洗、PIT 数据与确定性工具分置于三个进程；
        以 R0--R5 内环约束单次研判的可审计性，以 MI/LLR/HLR 三库记忆与浅/深反思
        外环约束跨时点一致性；以四层 PIT 过滤和内容寻址的证据绑定约束数据与引文。
  \item \textbf{评测}：设计并实现一套预注册的配对消融协议与评分工具链。协议固定
        压力/平静配对抽样、四档到二值的阈值、Wilson 区间、精确 McNemar、Holm 校正、
        覆盖率与遮蔽口径，并强制标签置换安慰剂和标签循环性诊断；工具链从真实产物
        生成论文表格与正文宏，缺失指标显示为 \texttt{--} 而非 0。
  \item \textbf{复现}：给出端到端复现命令、产物字段、模型标识、随机种子与提交
        占位；报告实现缺陷与负结果，使读者可以判断结论在何种条件下成立。
\end{enumerate}

\subsection{不主张的边界}

本文不主张：（i）多智能体协同必然提高预测准确率；（ii）观察到的差异具有因果
解释；（iii）系统已具备实时生产可用性——部分依赖序列已经停更，相关指标只能做
历史回测；（iv）结果可以外推到其他模型端点或提示词。上述边界在
第~\ref{sec:threats} 节逐项展开。

\section{相关工作}

\subsection{可验证计算与可审计日志}

可验证计算关注“如何让第三方相信一个计算结果确实由声明的程序在声明的输入上产生”。
Pinocchio 以简洁非交互证明把验证成本压到远低于重算\citep{parno2013pinocchio}。
本文不引入零知识证明，而采用更工程化的可重算路径：确定性工具在 Rust 侧产生
\texttt{inputs\_hash} 与 \texttt{evidence\_hash}，数值 claim 一律由工具结果覆盖，
核验层重新比对；这与“可验证”目标一致，但信任假设不同——它依赖本地二进制的
可复现构建与哈希口径，而非密码学证明。

安全审计日志研究强调日志的防篡改与可追溯\citep{schneier1999audit}。本文吸收其
“逐条绑定、不可静默改写”的原则：每次研判写入 \texttt{traces/<run\_id>.jsonl}，
证据以内容寻址 ID 绑定到 claim；核验失败或泄漏事件不会被删除，而是作为字段保留
并进入交付卡。与完整哈希链审计相比，本文的日志仍是单机、单写入者模型，未解决
恶意内核篡改问题。

\subsection{多智能体 LLM 协同与辩论}

多智能体辩论让多个模型实例交换答案并迭代修正，已被证明可改善事实性与推理
\citep{du2023debate}；另一类工作通过“鼓励分歧”的辩论协议提升发散思维
\citep{liang2023mad}。工程框架方面，AutoGen 把多智能体会话抽象为可编排的对话
\citep{wu2023autogen}，MetaGPT 以标准作业程序约束角色协作\citep{hong2023metagpt}。
这些工作为本文的三主体设计提供了直接参照，但它们的评估目标通常是答案质量或
任务完成度，而非“在固定时点上改变风险等级是否提升预警判别力”。本文的对照把
协同限制在证据共享、交叉验证、分歧分类与仲裁上，并用配对 McNemar 检验判别
差异，而不是比较生成文本的主观质量。

\subsection{检索增强、长期记忆与反思}

检索增强生成把外部知识引入解码过程\citep{lewis2020rag}；Generative Agents 用
记忆流与反思维持长期行为一致性\citep{park2023generative}；Reflexion 用语言化的
自我反馈替代梯度更新\citep{shinn2023reflexion}；MemGPT 把上下文管理类比为操作系统
的分页\citep{packer2023memgpt}。本文的 MI/LLR/HLR 三库记忆与这些工作共享“写入—
检索—反思”的结构，但增加两条领域约束：所有记忆条目携带 \texttt{available\_at}，
检索在向量库标量层强制 \texttt{available\_at <= as\_of}；反思回填使用结果可得日，
而不是事后复盘日。前者防止未来记忆污染历史研判，后者防止反思本身成为泄漏通道。

\subsection{金融压力测度、VaR 与 ES}

金融压力指数把多市场信息压缩为可比较的状态变量。Illing 与 Liu 为加拿大构建了
覆盖权益、债券、外汇与银行部门的压力指数\citep{illing2006measuring}；Holló、
Kremer 与 Lo Duca 提出的 CISS 用时变相关结构聚合五个子市场\citep{hollo2012ciss}；
Cardarelli、Elekdag 与 Lall 则把金融压力与实体经济收缩联系起来
\citep{cardarelli2011financial}。本文的 CFSI 是这一传统的六维实现，并用训练期
冻结分位数 $\tau$ 定义 \texttt{high\_stress}。

风险度量方面，VaR 的历史模拟与条件风险价值（CVaR/ES）是确定性工具链的基础。
Rockafellar 与 Uryasev 把 CVaR 转化为可优化的凸问题\citep{rockafellar2000cvar}；
Acerbi 与 Tasche 讨论 ES 的一致性\citep{acerbi2002coherence}。本文只用这些度量
作为证据生成器，不让 LLM 直接生成风险数值；数值 claim 由工具重算覆盖，这是
“零 Mock”红线的具体形式。

\subsection{预警评测方法与标签泄漏}

配对分类比较常用 McNemar 精确检验\citep{mcnemar1947note}；小样本比例区间用
Wilson 得分区间比正态近似更稳健\citep{wilson1927probable}；Dietterich 系统比较了
监督学习算法的近似统计检验\citep{dietterich1998approximate}。数据挖掘中的泄漏
研究指出，特征或标签构造中的时间错位会让离线指标系统性高估
\citep{kaufman2012leakage}；回测过拟合研究进一步说明，重复试验与选择性报告会
放大看似显著的收益\citep{bailey2016overfitting}。金融 LLM 评测方面，FinBen 提供
了覆盖多类金融任务的综合基准\citep{xie2024finben}，BloombergGPT 展示了领域预训练的
规模效应\citep{wu2023bloomberggpt}。本文不提出新的基准分数，而是把“配对抽样 +
时点过滤 + 循环性诊断 + 遮蔽统计”作为系统级消融的评测框架。

\subsection{本文定位}

与上述工作相比，本文的差异不在单点算法，而在把三类常被分开讨论的问题放到同一
条可复现流水线上：（i）多主体协同是否改变定级；（ii）改变是否以可审计、无泄漏的
方式发生；（iii）代价是多少。我们明确报告上一轮小样本对照的负结果，并把本轮
设计差异写进方法节，避免把“改造后的协同”与“改造前的协同”混为一谈。

\section{系统架构}

\subsection{三进程拓扑}

CBR-MAW 采用纯客户端三进程架构，不依赖中央服务器或常驻外部数据库服务。三个进程
通过标准输入输出的 NDJSON 传输 JSON-RPC 2.0 消息；Rust \texttt{kernel} 是工作区
存储的唯一写入者，切换工作区即重启内核。

\begin{table}[t]
\centering
\caption{三进程职责与硬边界}
\label{tab:processes}
\footnotesize
\begin{tabular}{>{\raggedright\arraybackslash}p{2.4cm}
                >{\raggedright\arraybackslash}p{5.4cm}
                >{\raggedright\arraybackslash}p{6.2cm}}
\toprule
进程 & 主要职责 & 硬边界 \\
\midrule
Dart 桌面外壳 & A2UI 交互、X-AlertAgent 编排、R0--R5、权限门卫、密钥、预算与审计流 &
不直接打开工作区存储；不重算内核哈希；无密钥时禁止 Mock \\
Go \texttt{collector} & 多源抓取、限流与备源、Inbox 清洗、五模态特征标准化、自更新 &
不承载风险定级；不写 PIT 分析表 \\
Rust \texttt{kernel} & DuckDB PIT 时序与面板、嵌入式 LanceDB 三库记忆、确定性数值工具、
黄金哈希与证据 ID & 单一写入者；所有查询强制 \texttt{available\_at <= as\_of} \\
\bottomrule
\end{tabular}
\end{table}

\subsection{双核嵌套：外环认知闭环与内环 R0--R5}

系统采用“双核嵌套”结构。外环 \texttt{X-AlertAgent} 负责跨时点认知：感知并归并
截至 \texttt{as\_of} 的五模态事件流，形成 MI 记忆；从 MI/LLR/HLR 三库检索并注入
当前上下文；执行浅层反思（短/中/长期风险变动归因与相似案例）与深层反思（回看
过去 $T$ 天的主体行动与事后预警效能）；最后输出预警三件套——风险值、四级等级
（由确定性阈值判定）与三主体离散行动及证据链方案。

内环 R0--R5 负责单次研判的可审计性。表~\ref{tab:r0r5} 给出各阶段的输入、动作与
产物。关键设计是：LLM 可以提出主张与解释，但不能直接改写工具数值；所有数值 claim
在物化阶段被工具结果覆盖，核验阶段再检查覆盖是否完整。

\begin{table}[t]
\centering
\caption{内环 R0--R5 的阶段、动作与失败处理}
\label{tab:r0r5}
\footnotesize
\setlength{\tabcolsep}{4pt}
\begin{tabular}{>{\raggedright\arraybackslash}p{1.1cm}
                >{\raggedright\arraybackslash}p{4.0cm}
                >{\raggedright\arraybackslash}p{4.2cm}
                >{\raggedright\arraybackslash}p{4.2cm}}
\toprule
阶段 & 动作 & 产物 & 失败/降级 \\
\midrule
R0 & 三主体在各自工具与信息白名单内独立分析，确定性规则给出初始等级 &
角色草稿、初始风险卡 & 信号不可用时按最低档占位并写限制说明 \\
R1 & 证据汇总到黑板，三级去重与别名解析 & 共享证据包 & 重复证据计数，不删除矛盾来源 \\
R2 & 工具确定性重算 + 同侪/裁判语义互评 & 复核结论 & 裁判不可用则标记 \texttt{available=false} \\
R3 & 分歧归一化为 fact / directional / preference / modality\_conflict &
分歧记录 & preference 单独计数，不算错误 \\
R4 & 按“工具重算 > 证据等级 > 人工”仲裁，结果进入决策损失 & 仲裁记录、降级状态 &
无法裁决时升级人工并保留双方证据 \\
R5 & 引文、PIT 泄漏、数值覆盖、权限、冲突、预算六类核验 + 语义核验 &
交付卡与 \texttt{limitations} & 泄漏事件 $>0$ 视为硬错误 \\
\bottomrule
\end{tabular}
\end{table}

\subsection{三库记忆：MI、LLR 与 HLR}

三库记忆使用嵌入式 LanceDB，物理上为单表 \texttt{memory}，以 \texttt{store}
判别列区分 \texttt{mi}、\texttt{llr}、\texttt{hlr}，差异字段放入 \texttt{payload}
JSON。这样一套检索代码即可完成跨库检索与 MMR 多样性重排。向量化使用
\texttt{fastembed} 与 BGE-small-zh-v1.5（512 维，离线确定性）；模型文件首次运行时
下载到工作区 \texttt{models/}，逐文件校验 SHA256，失败则显式降级为 FTS，而不是
伪造向量。

三个库的写入时机不同：MI 在感知阶段写入；LLR 在每次研判结束时写入；HLR 在写 LLR
时顺带回填过去 $T$ 天尚未结算的条目，其 \texttt{available\_at} 取“结果可得日”，
而不是“复盘执行日”。检索前在 LanceDB 标量层执行 \texttt{available\_at <= as\_of}
前过滤，再做 ANN/FTS 召回与 MMR 重排。该顺序保证过滤发生在候选集进入模型之前，
避免“先检索未来、再在提示词里删掉”的伪隔离。

\subsection{四层 PIT 过滤}

PIT 纪律由四层独立机制共同保证，任何一层失效都不能由其他层静默补偿：
\begin{enumerate}
  \item \textbf{采集层}：\texttt{collector} 为每条观测推导 \texttt{available\_at} 与
        \texttt{publish\_time}，区分“观测日”与“可得日”。
  \item \textbf{存储查询层}：DuckDB 与 LanceDB 的查询在 SQL/标量过滤中强制
        \texttt{available\_at <= as\_of}；禁止在应用层做“先取最新再裁剪”。
  \item \textbf{Agent 上下文层}：工具调用器强制注入 \texttt{as\_of}，禁止内核回落到
        当前时间；提示词只接收已过滤的证据包。
  \item \textbf{核验层}：R5 检查每条证据的时点，\texttt{leakage\_events > 0} 为硬错误；
        交付卡记录泄漏计数，而不是只记录“通过”。
\end{enumerate}
这四层对应实现中曾经出现的真实缺陷：工具调用未传 \texttt{as\_of} 会让历史时点读到
“今天”的数据，而应用层仍按请求时点盖章，泄漏检测器因此永远看不到问题。把
\texttt{as\_of} 注入下沉到工具调用器，是该缺陷的修复方式。

\subsection{证据哈希绑定与核验}

证据以内容寻址方式绑定：工具名、版本、\texttt{as\_of} 与参数共同决定
\texttt{inputs\_hash} 与 \texttt{evidence\_id}；工具输出连同输入哈希生成
\texttt{evidence\_hash}。Rust 内核是哈希口径的唯一权威，Dart 侧只消费内核返回值，
不重算三级哈希，避免两套浮点规范分叉。每条 claim 记录 \texttt{evidence\_ids}；
物化阶段用工具值覆盖 LLM 自报数值，并单独计数“LLM 自报值 vs 工具重算值”的不一致，
避免“先覆盖再比对必然全过”的假阴性。

预算与审计使用 $J = L + \mu C$ 记录决策损失与计算成本；每步写入
\texttt{traces/<run\_id>.jsonl}。预算耗尽时标记 \texttt{partial=true} 并输出已验证
部分结论，而不是静默截断。

\subsection{三主体与离散行动}

三主体分别为监管（\texttt{regulator}）、机构（\texttt{institution}）与投资者
（\texttt{investor}）。监管的行动空间为 \texttt{In}、\texttt{Watch}、
\texttt{Out}，机构为 \texttt{Discount}、\texttt{Hedge}、\texttt{Equilibrium}，
投资者为 \texttt{Buy}、\texttt{Hold}、\texttt{Sell}。另有专职交叉验证角色
\texttt{crossValidator}，只输出 \texttt{confirmed}/\texttt{refuted}/\texttt{uncertain}，
不产出主风险卡。三主体的 LLM 自评等级与确定性定级分开记录；两者不一致时，
系统只允许自评“上调”等级（漏报成本高于误报），下调一律不采纳，并在
\texttt{limitations} 中写明调整原因。

\subsection{消融开关与实验锁定}

同一套主代码通过统一开关控制协作模式（\texttt{single}/\texttt{independent}/
\texttt{collaborative}）、记忆库（MI/LLR/HLR）、多模态辅助与证据强度。跑论文对照时
设置 \texttt{experiment.lock\_models=true}，绕过能力槽位路由，把所有生成槽位锁定到
同一模型端点，消除模型异质性这一混杂变量。基准控制台默认开启记忆（除非显式传
\texttt{--no-memory}），三臂共享同一工具注册表、同一信号源与同一模型；差异只来自
协作模式。

\section{实验协议}
\label{sec:protocol}

\subsection{数据与标签}

实验数据来自真实 PIT 库（60313 条观测、27 个序列、9798 篇文档），日历覆盖
2015-01-05 至 2026-10-28，共 2902 个交易日。标签为 \texttt{stress\_pca\_q90}：
在训练期（\texttt{train\_end=2021-12-31}）冻结的分位数 $\tau$ 上定义
\texttt{high\_stress}，再用未来 20 个交易日的压力状态构造二值标签
$y_{20,t}=\max(hs_{t+1},\dots,hs_{t+20})$，即严格未来、不含当日。
训练期之后为样本外评估区间。本文对随附黄金向量逐行核验了该口径
（$y_{20}$ 与 $t+1$ 至 $t+20$ 窗口在 \RESULT{label}{n_agree}/
\RESULT{label}{n_future} 个可定义日上一致），并另设更严格标签
$y^{\mathrm{strict}}_{t}=\max(hs_{t+2},\dots,hs_{t+21})$ 用于量化压力状态自相关。

标签与指标同源是本协议最需要正面处理的问题：\texttt{high\_stress} 本身由 CFSI
超过 $\tau$ 定义，而任何以 CFSI 为输入的臂都在“重建标签”。本文不掩盖这一点，
而是在结果节给出条件概率 $P(y_{20}\mid hs_t)$ 与更严格标签的诊断，并要求读者
按折扣解读 lift。项目内部报告曾把 $y_{20}$ 写成“含当日”；逐行核验表明该表述错误，
现已在文档与生成器中统一修正为严格未来口径，核验过程与差异记录于附录~\ref{app:repro}。

\subsection{配对抽样}

LLM 在环臂的每次研判有真实成本，不能对整段日历穷举。协议采用配对抽样：
\begin{enumerate}
  \item 只取 $y_{20}$ 非空的时点（末尾 20 行天然无标签，若纳入会把所有臂
        都算成负例）；
  \item 压力期与平静期各取一半，避免“只测平静期所有臂都正确”或“只测压力期
        所有臂都过警”；
  \item 两组内部按时间均匀抽取，使两组的日历分布相近，降低“压力样本集中在某一年”
        的混杂；
  \item 固定种子（本稿产物 \texttt{seed=\RESULT{meta}{seed}}），抽样结果写入产物，
        合并分片时拒绝种子或模型不一致的输入。
\end{enumerate}
正式实验为 120 个配对时点（压力/平静各 60），分为 4 个分片并行执行；分片只切分
同一份抽样，合并后与单进程等价。合并脚本对同一 \texttt{as\_of} 出现两次直接报错，
不静默取其一。

\subsection{三个对照臂}

\begin{table}[t]
\centering
\caption{三臂定义（同一模型、同一工具注册表、同一记忆配置）}
\label{tab:arms}
\footnotesize
\begin{tabular}{>{\raggedright\arraybackslash}p{2.9cm}
                >{\raggedright\arraybackslash}p{11.2cm}}
\toprule
臂 & 定义 \\
\midrule
\texttt{collaborative} & 完整 R0--R5：三主体独立分析 $\rightarrow$ 证据黑板共享
$\rightarrow$ 交叉验证 $\rightarrow$ 分歧分类 $\rightarrow$ 仲裁 $\rightarrow$
六类核验；包含跨角色证据共享与冲突处理。 \\
\texttt{single} & 一个通用 agent 依次完成三主体任务；有工具与证据链，但不做
跨角色证据共享、交叉验证与仲裁。 \\
\texttt{independent} & 三主体各自独立决策；不共享证据、不交叉验证、不协调，
保留各自的工具与定级流程。 \\
\bottomrule
\end{tabular}
\end{table}

三臂使用同一生成端点与同一模型（\texttt{\RESULT{meta}{model}}），
\texttt{lock\_models=true}；同一工具集、同一信号源、同一 PIT 过滤与同一预算上限。
基准控制台在每次运行时创建独立工作区，但同一分片内三臂按固定顺序执行并共享记忆库；
这一潜在的同日跨臂检索污染在威胁节讨论。

\subsection{二值化与指标}

四档等级为 \texttt{low}/\texttt{medium}/\texttt{high}/\texttt{critical}。主分析把
\texttt{high} 及以上视为“预警”（阈值 \texttt{high}），与协议冻结口径一致；
\texttt{critical} 作为敏感性阈值可由脚本参数切换。给定预测 $p_t\in\{0,1\}$ 与
真值 $y_t\in\{0,1\}$：
\begin{gather}
  \mathrm{precision} = \frac{TP}{TP+FP}, \qquad
  \mathrm{recall} = \frac{TP}{TP+FN}, \\
  \mathrm{lift} = \frac{\mathrm{precision}}{\text{基准发生率}}, \qquad
  \text{触发率} = \frac{TP+FP}{n}.
\end{gather}
基准发生率为同一批配对时点中正例的占比。均衡准确率为正负例召回率的平均：
$(\mathrm{TPR}+\mathrm{TNR})/2$，在压力/平静各半的抽样下比 precision 更能区分
“会不会区分”与“触发多少”。同时报告平均端到端延迟（秒）与平均 token 消耗。

\subsection{统计推断}

precision 与 recall 的 95\% 区间使用 Wilson 得分区间\citep{wilson1927probable}；
配对比较使用双侧精确 McNemar 检验\citep{mcnemar1947note}，不一致对数为 0 时
报告“无不一致对”而不报 $p=1$。两项比较（\texttt{single} 与
\texttt{independent} 分别对 \texttt{collaborative}）使用 Holm 校正。按协议，
$n<30$ 或无不一致对时只报描述统计，不作显著性主张；区间宽度本身也是结果的一部分，
不能因为区间宽就省略。

\subsection{安慰剂、循环性与遮蔽}

\paragraph{标签置换安慰剂。} 保持每个臂的预测向量与正例数不变，随机置换真值顺序
$B$ 次，重算 lift，得到零分布。若观测 lift 落在零分布的常见范围内，说明该向量
相对随机标签分配并不异常。该检验只能在评分向量层面回答“是否异常”，不能替代
独立重跑 LLM 的安慰剂；两者在结果节分开标注。

\paragraph{标签循环性诊断。} 报告 $P(y_{20}=1\mid hs_t=1)$、
$P(y_{20}=1\mid hs_t=0)$，以及更严格标签
$y^{\mathrm{strict}}_t=\max(hs_{t+2},\dots,hs_{t+21})$ 下的同组条件概率。
若 $y_{20}$ 与 $t+1$ 至 $t+20$ 窗口在全部可定义日上一致，说明本产物的 $y_{20}$
不含当日；此时循环性来自 $hs_t$ 与 CFSI 的同源关系，以及压力状态的自相关。
更严格标签把起点再推后一天，条件概率的下降幅度就是自相关折扣的量化。

\paragraph{覆盖率与遮蔽。} 覆盖率 $=$ 有效定级数 / 计划评分时点数；遮蔽表分别
统计无有效定级、运行失败、降级角色数、泄漏事件、无证据主张与核验未通过。遮蔽
记录不会被改写成负例；缺失指标在表中显示为 \texttt{--}，不以 0 代替。

\subsection{预注册、偏离与多重比较}

协议在正式跑批前固定：抽样种子、压力/平静各半、阈值 \texttt{high}、两项配对比较、
Holm 校正、$n<30$ 不作显著性主张、缺失不补数。任何偏离（例如端点变更、提示词
修改、抽样重跑）都必须记录在附录~\ref{app:repro} 的偏离表中；若偏离影响口径，
结果节必须并列报告偏离前后两组数字，而不是只报较好看的一组。

\section{结果}
\label{sec:results}

本节所有数字由 \texttt{tools/bench/analysis/render\_tables.py} 从
\texttt{artifacts/bench} 的真实产物生成；正文通过 \verb|\RESULT{臂}{指标}| 宏取值，
表格由同一脚本写入 \texttt{tables/}。当前产物：
由 \RESULT{meta}{n_shards} 个分片合并，记录 \RESULT{meta}{n_records} 条，计划评分时点
\RESULT{meta}{n_planned} 个，已进入评分的时点 \RESULT{meta}{n_scored} 个，
产物是否仍为部分结果：\RESULT{meta}{partial}。若某单元格显示 \texttt{--}，
表示该指标在产物中不可用，而不是取值为 0。

跑批过程中修正了一处编排预算配置：旧批次把每角色 token 上限设得低于真实消耗，
导致每次 \texttt{chargeTokens} 都截断、\texttt{partial} 恒为真。旧批次已全部
丢弃，本稿表格只使用修正预算后的新批次产物。编译时产物是否为部分结果：
\RESULT{meta}{partial}。若仍为“是”，表格只是中间快照，不能作为最终效应量引用；
最终数字应在跑批完成后由同一脚本重新生成。

\subsection{主表：判别力、成本与触发率}

表~\ref{tab:main} 给出三臂在同一批配对时点上的 precision、Wilson 区间、recall、
lift、触发率、均衡准确率、平均延迟与平均 token。$\mathrm{lift}>1$ 表示 precision
高于基准发生率；均衡准确率用于在压力/平静各半的抽样下比较“区分能力”而非“触发
倾向”。

\begin{table}[t]
\centering
\caption{三臂预警判别结果（二值化阈值：\texttt{high}；缺指标显示 \texttt{--}）}
\label{tab:main}
\footnotesize
\setlength{\tabcolsep}{3pt}
\begin{tabular}{lrrrrrrrrr}
\toprule
臂 & $n$ & precision & 95\% CI & recall & lift & 触发率 & 均衡准确率 & 耗时(s) & token \\
\midrule
协同 C & \RESULT{collaborative}{n} & \RESULT{collaborative}{precision} &
\RESULT{collaborative}{precision_ci} & \RESULT{collaborative}{recall} &
\RESULT{collaborative}{lift} & \RESULT{collaborative}{trigger} &
\RESULT{collaborative}{balacc} & \RESULT{collaborative}{latency_s} &
\RESULT{collaborative}{tokens} \\
单智能体 S & \RESULT{single}{n} & \RESULT{single}{precision} &
\RESULT{single}{precision_ci} & \RESULT{single}{recall} &
\RESULT{single}{lift} & \RESULT{single}{trigger} &
\RESULT{single}{balacc} & \RESULT{single}{latency_s} &
\RESULT{single}{tokens} \\
独立 I & \RESULT{independent}{n} & \RESULT{independent}{precision} &
\RESULT{independent}{precision_ci} & \RESULT{independent}{recall} &
\RESULT{independent}{lift} & \RESULT{independent}{trigger} &
\RESULT{independent}{balacc} & \RESULT{independent}{latency_s} &
\RESULT{independent}{tokens} \\
\bottomrule
\end{tabular}
\par\smallskip
\raggedright\scriptsize
注：$n$ 为进入评分的配对数；CI 为 precision 的 Wilson 95\% 区间；
lift $=$ precision $/$ 基准发生率；耗时与 token 为各臂在全部已跑记录上的均值。
缺失指标（含无正例预测时的 precision）写作 \texttt{--}，不以 0 代替。
若任一臂 $n<30$，按协议不作显著性主张。
\end{table}

对 H1 的方向性判断不能只看点估计：协同臂与单臂的 precision 差为
\RESULT{collaborative}{precision} 对 \RESULT{single}{precision}，但区间宽度与配对
结构决定该差是否可解释。对 H2，成本差异见同表后两列：协同臂平均延迟
\RESULT{collaborative}{latency_s} 秒、平均 token \RESULT{collaborative}{tokens}，
单智能体臂为 \RESULT{single}{latency_s} 秒与 \RESULT{single}{tokens}。H1/H2 的
正式结论：协同臂相对单智能体臂「\RESULT{concl}{vs_single}」、相对独立臂
「\RESULT{concl}{vs_indep}」（Holm 校正后）；成本为 token \RESULT{concl}{token_ratio} 倍、
耗时 \RESULT{concl}{latency_ratio} 倍。

\subsection{配对 McNemar 检验}

表~\ref{tab:mcnemar} 报告两项配对比较。$b$ 为“仅协同正确”的不一致对，$c$ 为
“仅基线正确”的不一致对；$p$ 为双侧精确二项检验，Holm $p$ 覆盖两项比较。

\begin{table}[t]
\centering
\caption{配对 McNemar 精确检验（以协同臂为处理组）}
\label{tab:mcnemar}
\footnotesize
\setlength{\tabcolsep}{4pt}
\begin{tabular}{lrrrrrr}
\toprule
比较 & $n$ & 仅协同正确 & 仅基线正确 & $p$ & Holm $p$ & 判定 \\
\midrule
单智能体 S vs 协同 C & \RESULT{mcnemar_single}{n} &
\RESULT{mcnemar_single}{only_collab} & \RESULT{mcnemar_single}{only_base} &
\RESULT{mcnemar_single}{p} & \RESULT{mcnemar_single}{p_holm} &
\RESULT{mcnemar_single}{verdict} \\
独立 I vs 协同 C & \RESULT{mcnemar_independent}{n} &
\RESULT{mcnemar_independent}{only_collab} & \RESULT{mcnemar_independent}{only_base} &
\RESULT{mcnemar_independent}{p} & \RESULT{mcnemar_independent}{p_holm} &
\RESULT{mcnemar_independent}{verdict} \\
\bottomrule
\end{tabular}
\par\smallskip
\raggedright\scriptsize
注：双侧精确二项检验；Holm 校正覆盖两项比较。$n<30$ 或无不一致对时只报描述统计，
不作显著性主张。当前判定：单智能体 \RESULT{mcnemar_single}{verdict}；
独立 \RESULT{mcnemar_independent}{verdict}。
\end{table}

若两项检验均不显著，正确表述是“在本样本与本案端点下，没有检测到协同带来的判别
差异”，而不是“协同无效”。若其中一项显著，仍需检查方向是否与 H1 一致、是否
稳健于标签循环性诊断与安慰剂零分布；单次显著不作为强主张。

\subsection{覆盖率与遮蔽}

表~\ref{tab:coverage} 与表~\ref{tab:masking} 把“能不能给出定级”和“给得是否可核验”
分开报告。覆盖率低不一定意味着模型差：它可能来自预算耗尽、端点失败、分片未完成
或记录尚未合并。遮蔽计数用于区分这些原因。

\begin{table}[t]
\centering
\caption{覆盖率：有效定级占计划评分时点的比例}
\label{tab:coverage}
\footnotesize
\setlength{\tabcolsep}{3.5pt}
\begin{tabular}{lrrrrrrr}
\toprule
臂 & 计划时点 & 已跑记录 & 有效定级 & 覆盖率 & partial & 有降级 & 核验通过率 \\
\midrule
协同 C & \RESULT{collaborative}{planned} & \RESULT{collaborative}{attempted} &
\RESULT{collaborative}{valid} & \RESULT{collaborative}{coverage} &
\RESULT{collaborative}{partial_rate} & \RESULT{collaborative}{degraded_rate} &
\RESULT{collaborative}{verify_pass} \\
单智能体 S & \RESULT{single}{planned} & \RESULT{single}{attempted} &
\RESULT{single}{valid} & \RESULT{single}{coverage} &
\RESULT{single}{partial_rate} & \RESULT{single}{degraded_rate} &
\RESULT{single}{verify_pass} \\
独立 I & \RESULT{independent}{planned} & \RESULT{independent}{attempted} &
\RESULT{independent}{valid} & \RESULT{independent}{coverage} &
\RESULT{independent}{partial_rate} & \RESULT{independent}{degraded_rate} &
\RESULT{independent}{verify_pass} \\
\bottomrule
\end{tabular}
\par\smallskip
\raggedright\scriptsize
注：覆盖率 $=$ 有效定级 / 计划评分时点；partial 与“有降级”为占已跑记录的比例，
其中 partial 是运行摘要字段的原值（预算/时间截断标记），不是本文的判定；
核验通过率为 verification 非空记录中的通过比例。缺失指标显示 \texttt{--}。
\end{table}

\begin{table}[t]
\centering
\caption{遮蔽计数：被排除在评分之外或未通过核验的记录}
\label{tab:masking}
\footnotesize
\setlength{\tabcolsep}{3.5pt}
\begin{tabular}{lrrrrrr}
\toprule
臂 & 无有效定级 & 运行失败 & 降级角色数 & 泄漏事件 & 无证据主张 & 核验未通过 \\
\midrule
协同 C & \RESULT{collaborative}{null_risk} & \RESULT{collaborative}{failures} &
\RESULT{collaborative}{degraded_roles} & \RESULT{collaborative}{leakage} &
\RESULT{collaborative}{unsupported} & \RESULT{collaborative}{verify_failed} \\
单智能体 S & \RESULT{single}{null_risk} & \RESULT{single}{failures} &
\RESULT{single}{degraded_roles} & \RESULT{single}{leakage} &
\RESULT{single}{unsupported} & \RESULT{single}{verify_failed} \\
独立 I & \RESULT{independent}{null_risk} & \RESULT{independent}{failures} &
\RESULT{independent}{degraded_roles} & \RESULT{independent}{leakage} &
\RESULT{independent}{unsupported} & \RESULT{independent}{verify_failed} \\
\bottomrule
\end{tabular}
\par\smallskip
\raggedright\scriptsize
注：无有效定级指 risk 缺失或不在四档枚举内；泄漏事件与无证据主张为 verification
字段的累计计数；这些记录不会被改写成负例。
\end{table}

覆盖率与判别力必须一起读：若协同臂覆盖率显著低于单臂，任何“协同更准”的比较都
受选择偏差影响；若协同臂覆盖率更高但判别力不变，说明协同改善的是“出结论的稳定性”
而非“区分能力”。若覆盖率表中的 partial 列接近 1，说明所有已跑记录都带部分结论
标记；此时主表数字仍应报告，但不能表述为“完整预算下的协同效应”，必须先确认标记
来自预算/时间截断还是字段语义。当前产物下的关系：最小覆盖率 \RESULT{concl}{min_coverage}，最高降级率
\RESULT{concl}{max_degraded}，部分结论标记（\RESULT{concl}{partial_any}）。
覆盖率不可比时，主表结论只能表述为「在已有产物上的判别力」，不能外推为完整预算下的效应。

\subsection{安慰剂与标签循环性}

表~\ref{tab:placebo} 给出评分向量层面的标签置换安慰剂；表~\ref{tab:circularity}
给出标签循环性诊断。前者回答“预测向量相对随机标签是否异常”，后者回答“标签与
指标同源带来多少折扣”。

\begin{table}[t]
\centering
\caption{标签置换安慰剂：打乱真值顺序后的 lift 零分布}
\label{tab:placebo}
\footnotesize
\setlength{\tabcolsep}{4pt}
\begin{tabular}{lrrrrrr}
\toprule
臂 & 观测 lift & 零分布均值 & 零分布 95\% 分位 & $p$ & $B$ & 判定 \\
\midrule
协同 C & \RESULT{placebo_collaborative}{obs_lift} &
\RESULT{placebo_collaborative}{null_mean} & \RESULT{placebo_collaborative}{null_q95} &
\RESULT{placebo_collaborative}{p} & \RESULT{placebo_collaborative}{draws} &
\RESULT{placebo_collaborative}{verdict} \\
单智能体 S & \RESULT{placebo_single}{obs_lift} &
\RESULT{placebo_single}{null_mean} & \RESULT{placebo_single}{null_q95} &
\RESULT{placebo_single}{p} & \RESULT{placebo_single}{draws} &
\RESULT{placebo_single}{verdict} \\
独立 I & \RESULT{placebo_independent}{obs_lift} &
\RESULT{placebo_independent}{null_mean} & \RESULT{placebo_independent}{null_q95} &
\RESULT{placebo_independent}{p} & \RESULT{placebo_independent}{draws} &
\RESULT{placebo_independent}{verdict} \\
\bottomrule
\end{tabular}
\par\smallskip
\raggedright\scriptsize
注：保持预测向量与正例数不变，仅置换真值顺序；
$p=(1+\#\{\mathrm{lift}_{\mathrm{null}}\ge \mathrm{lift}_{\mathrm{obs}}\})/(B+1)$。
该检验不替代独立重跑的安慰剂。
\end{table}

\begin{table}[t]
\centering
\caption{标签循环性诊断（$y_{20}$ 的未来窗口与更严格标签）}
\label{tab:circularity}
\footnotesize
\begin{tabular}{lr}
\toprule
统计量 & 值 \\
\midrule
样本日数（$y_{20}$ 可用） & \RESULT{label}{n} \\
基准发生率 $P(y_{20}=1)$ & \RESULT{label}{base_y20} \\
$P(y_{20}=1\mid hs_t=1)$ & \RESULT{label}{p_y_hs1} \\
$P(y_{20}=1\mid hs_t=0)$ & \RESULT{label}{p_y_hs0} \\
$t+1$ 至 $t+20$ 未来窗口样本日数 & \RESULT{label}{n_future} \\
$y_{20}$ 与 $t+1$ 至 $t+20$ 窗口一致的日数 & \RESULT{label}{n_agree} \\
更严格标签（$t+2$ 至 $t+21$）样本日数 & \RESULT{label}{n_strict} \\
基准发生率 $P(y^{\mathrm{strict}}=1)$ & \RESULT{label}{base_strict} \\
$P(y^{\mathrm{strict}}=1\mid hs_t=1)$ & \RESULT{label}{p_strict_hs1} \\
$P(y^{\mathrm{strict}}=1\mid hs_t=0)$ & \RESULT{label}{p_strict_hs0} \\
\bottomrule
\end{tabular}
\par\smallskip
\raggedright\scriptsize
注：$hs_t$ 为当日 \texttt{high\_stress}（定义即 CFSI $>\tau$）。条件概率只能作为
循环性诊断，不能当作预测能力。$y_{20}$ 与 $t+1$ 至 $t+20$ 窗口完全一致，说明
$y_{20}$ 不含当日；更严格标签的下降幅度用于量化自相关折扣。
\end{table}

解读规则：若 $P(y_{20}=1\mid hs_t=1)$ 明显高于基准发生率，说明“当日在压力区”
与“未来 20 日进入压力区”高度重合，任何以 CFSI 为输入的臂都会获得部分定义性
优势；此时 lift 必须与更严格标签下的 lift 并列报告，而不能只报较大的一个。
$y_{20}$ 与 $t+1$ 至 $t+20$ 窗口完全一致，说明本产物的 $y_{20}$ 不含当日；
折扣应归因于 $hs_t$ 与 CFSI 的同源关系以及压力状态自相关，而不是“含当日”的
机械重合。

\subsection{负结果与异常}

我们要求结果节显式记录以下情形，而不是把它们移到脚注：
\begin{itemize}
  \item 任一臂的 $n<30$、不一致对为 0、或置信区间宽到无法区分实际差异；
  \item 协同臂的覆盖率或核验通过率低于单臂，导致判别比较不可比；
  \item 安慰剂 $p$ 不小，说明观测 lift 与随机标签分配不可区分；
  \item 更严格标签下的条件概率与 $y_{20}$ 差异很大，且折扣幅度改变结论方向；
  \item 运行失败、降级或分片缺失导致某些时点未被评分；
  \item \texttt{partial} 标记占比异常（例如全部记录为 true），且无法用预算或时间
        上限解释。
\end{itemize}
当前产物下的具体异常：最小覆盖率 \RESULT{concl}{min_coverage}（$<1$ 即有时点未评分），
最高降级率 \RESULT{concl}{max_degraded}，部分结论标记（\RESULT{concl}{partial_any}）。
泄漏事件数、核验失败数与不支持主张数分别见覆盖率表三列；任一非零即计入本节异常清单。

\subsection{记忆消融（MI/LLR/HLR 开与关）}
\label{sec:memory-ablation}

为回答「三库记忆是否改变判别」，在同一端点、同一提示词与同一工具集下，对
\texttt{collaborative} 臂做一次开/关对照：记忆关组把记忆客户端置空，\emph{三库记忆}（MI 检索、LLR/HLR 写回）都不发生
（不是「忽略结果」）；两条条件都保留 DuckDB 侧的相似案例反思，因此本对照隔离的是
LanceDB 三库记忆本身，而非「有没有反思」。操作检查先行：\texttt{memory\_written} 事件在
记忆开组出现、记忆关组为 0，确认开关确实生效。

配对时点 \RESULT{mem}{n} 个。均衡准确率：记忆开 \RESULT{mem}{balacc_a}、
记忆关 \RESULT{mem}{balacc_b}，差 \RESULT{mem}{balacc_delta}
（配对 McNemar $p=\RESULT{mem}{p}$，判定：\RESULT{mem}{verdict}）。

读法与主表一致：若两条件预测一致，正确结论是「记忆\emph{不改变判别}，其价值在
可解释性与跨时点归因」，而\emph{不是}「记忆无效」——后者需要变更判别的证据，
本设计并未提供。

\section{讨论}

\subsection{上一轮负结果}

上一轮在 12 个配对时点、四臂（含纯 LLM 直答）上实跑，得到本文最重要的工程教训：
\texttt{single}、\texttt{independent} 与 \texttt{collaborative} 三臂的预测向量
逐点一致，precision/recall/lift 完全相同（lift 均为 1.09）；协同臂多花约 20 秒与
6k token，换来同一组预测\footnote{见 \texttt{docs/13-基准测试报告.md} §3.5.1。
该轮 $n=12$，低于本协议的 30 例门槛，只作工程归因，不作显著性主张。}。
代码级归因有两条：定级由确定性规则给出，三角色的 LLM 自评只记录、不参与定级；
三主体共享同一信号源与同一工具计划，文本解读的差异无法进入等级。

这不是“协同没用”的终局结论，而是“当时的协同没有把多主体信息接入决策”的实现
事实。它直接决定了本轮的设计差异：如果自评仍然不进入定级，那么增加角色、增加
辩论轮次都只是在同一份确定性信号上生成更多解释，判别向量不会改变。

\subsection{本轮差异设计}

本轮针对上述归因做了三项改造，并在方法节明确记录，避免把两轮结果混为一谈：
\begin{enumerate}
  \item \textbf{LLM 自评非对称进入定级}：当自评等级高于确定性定级时上调一级，
        低于时一律不采纳。该设计把模型当作“额外报警源”，利用漏报成本高于误报的
        决策损失权重，同时不允许一次幻觉压低真实风险。
  \item \textbf{三主体分歧跨度}：记录三主体自评等级的极差
        \texttt{assessmentSpread}；单智能体恒为 0，协同臂可大于 0。它把“角色之间
        是否真的不同意”变成可观测的过程量，而不是只看最终等级。
  \item \textbf{三主体离散行动}：监管、机构、投资者各自输出枚举行动及证据链，
        使协同的产物不止一个等级，还包括可审计的行动建议。
\end{enumerate}
这三项改造是否足以让三臂在判别指标上分开，是 H1 的实证问题。若仍然没有差异，
可能的机制包括：自评上调只在少数时点触发；三主体共享同一信号源，分歧集中在
文本而非等级；或者预警判别力主要由 CFSI 映射决定，协同的解释层无法改变阈值附近
的排序。讨论这些机制时，必须用分歧跨度、降级计数与遮蔽表作为证据，而不是凭
印象断言“模型在讨论”。

\subsection{成本、可审计性与产品含义}

即使 H1 被证伪，协同仍可能改变三类非判别指标：证据可溯源率、数值不一致检出与
冲突分类覆盖。这些指标不直接提高 precision，但它们决定系统能否被监管、机构与
投资者信任。反过来，如果协同只增加成本而不增加判别力与可审计性，那么在产品上
应默认关闭协同，把预算留给数据新鲜度、PIT 纪律与工具覆盖。

本文的立场是“可审计性是独立目标，不以判别力提升为前提”。这也是为什么结果节
同时报告覆盖率、遮蔽与核验通过率：一个精度略高但大量记录被降级或核验失败的臂，
不应被当作更好的预警系统。

\subsection{对后续实验的设计含义}

本轮协议留下了三个可复用的设计约束。第一，配对抽样必须在压力/平静两层内进行，
否则“所有臂都正确”或“所有臂都过警”会让对照失去分辨力。第二，标签与指标同源
时，必须把循环性诊断写进主结果，而不是放在附录；否则 lift 会被读者误读为预测
能力。第三，LLM 在环实验的成本决定了样本量，样本量又决定了可主张的效应量；
在 $n$ 较小时，报告区间与“不作显著性主张”比追求 $p<0.05$ 更诚实。

\section{有效性威胁}
\label{sec:threats}

\subsection{未运行的消融维度}

本文只报告了协作模式这一维度的正式对照（三臂 × 120 配对时点）。以下开关虽已实现
并可运行，但\textbf{未纳入本轮正式实验}，因此不给出任何效应量：

\begin{itemize}
  \item \textbf{三库记忆开/关}（\texttt{--no-memory}）：机制已接线并跑通
        （MI/LLR 写入与检索、HLR 结算），但未做系统性的开/关对照；
  \item \textbf{能力槽位异构路由}（ADR-011）：本轮锁定单一端点，无可比对象；
  \item \textbf{弱证据组}（\texttt{evidence\_strength=weak}）：用于测幻觉率，
        本轮未跑；
  \item \textbf{模态子集}：语音/图像/视频解析未实现，五模态只有文本与数值。
\end{itemize}

把这些维度说成“已验证无效果”是最容易犯的错，本文明确避免：未跑即未跑。

\subsection{标签循环性}

\texttt{high\_stress} 由 CFSI 超过 $\tau$ 定义，而 $y_{20}$ 由未来的
\texttt{high\_stress} 构造（$t+1$ 至 $t+20$，不含当日）；任何以 CFSI 为输入的臂
都在部分重建标签。本文用 $P(y_{20}\mid hs_t)$ 与更严格标签量化该问题，但无法完全
消除：即使更严格标签下 lift 仍大于 1，也只能说明信号并非完全循环，不能说明它对
真实风险有独立信息。
更根本的检验需要外生标签（例如事后监管处置、机构违约或市场流动性事件），本文
数据不包含这些标签。

\subsection{样本量与检验功效}

正式实验为 120 个配对时点，但 McNemar 只使用不一致对；当三臂预测高度一致时，
不一致对很少，检验功效很低。此时“不显著”不能解释为“无差异”，只能解释为“本
样本无法区分”。同样，覆盖率、降级与遮蔽是计数指标，在小样本下波动很大；本文
只作描述性报告，不把它们包装成显著结论。

\subsection{记忆库共用的历史缺陷（已修复）}

主实验跑批期间，记忆库目录按工作区父目录解析，导致同目录下的多个分片共用一份
记忆库；该缺陷在记忆消融的操作检查中被发现（全新工作区首次检索即命中历史记录），
随后改为按工作区隔离。影响范围：检索仍强制时点前过滤，故不构成前视偏差；同一时点
上三臂共享同一份污染，配对检验结论不受影响；但绝对 lift 含额外噪声，应视为上界。
正式投稿前应以隔离版重跑主表。

\subsection{单一模型端点与非确定性}

所有臂锁定到同一模型端点（\texttt{\RESULT{meta}{model}}），这消除了模型异质性，
但也把结论限制在该端点与当时的提示词版本上。LLM 生成具有非确定性：同一输入重跑
可能得到不同等级，尤其是在阈值附近。本文没有对每个时点重复采样，因此无法把
“模型随机性”与“协作模式差异”完全分离。跨模型外推需要独立实验。

\subsection{数据停更与实时性}

部分依赖序列已经停更：北向资金最新到 2024-08-16，南向、信用利差、回购与
USDCNY 等到 2025-09-30。依赖这些序列的工具在停更后只能做历史回测，不能支撑
“实时预警”的叙事。本文的评估区间包含停更前后，若某臂在停更后触发率骤降，必须
区分“信号失效”与“数据缺失”；遮蔽表与新鲜度检查用于此目的，但本文不把停更
后的结果外推为生产能力。

\subsection{其他威胁}

\begin{itemize}
  \item \textbf{提示词与实现漂移}：三臂共享同一提示词与工具注册表，但协同臂的
        编排路径更长，任何路径特有的 bug 都可能被误读为“协同效应”。上一轮已经
        出现过“工具名被当风险关键词”和“A 组面板为空”两类 bug；本轮要求把
        实现缺陷写入复现附录。
  \item \textbf{同日跨臂记忆污染}：同一分片内三臂按固定顺序执行并共享记忆库；
        若写入的 \texttt{available\_at} 不晚于当日，后执行的臂可能检索到先执行臂
        写入的记忆。该风险依赖记忆写入时点，本文未做逐条审计，属于需要后续隔离
        实验确认的混淆。
  \item \textbf{选择性抽样}：配对抽样固定种子，且只取有标签的时点；但日历分布
        仍可能与真实使用场景不同，结果不应外推到未抽样年份。
  \item \textbf{无人类评估}：本文没有让监管或机构从业者对行动建议做盲评；
        “可审计”是结构属性，不等于“可被采纳”。
  \item \textbf{基准污染}：模型训练数据可能包含公开市场历史；本文无法排除模型
        对特定历史时期的记忆，这会影响所有 LLM 在环臂，而不只是协同臂。
\end{itemize}

\section{结论}

本文提出并实现 CBR-MAW，一个以可审计性、时点一致性与可复现消融为核心约束的
跨境金融风险多主体预警系统；并给出一套从配对抽样到表格生成的评测流水线。本文
不预设协同一定带来判别增量：预注册假设、Wilson 区间、精确 McNemar、Holm 校正、
标签置换安慰剂与循环性诊断共同构成证伪条件。当前实验结论：
协同臂相对单智能体臂「\RESULT{concl}{vs_single}」、相对独立臂「\RESULT{concl}{vs_indep}」，
均衡准确率差分别为 \RESULT{concl}{balacc_delta_single} 与
\RESULT{concl}{balacc_delta_indep}，预测向量一致性：\RESULT{concl}{identical}；
代价是 token \RESULT{concl}{token_ratio} 倍、耗时 \RESULT{concl}{latency_ratio} 倍。
该判定仅在覆盖率 \RESULT{concl}{min_coverage} 的已评分时点上成立。
无论结果方向如何，本文的工程贡献与负结果记录都为后续“多主体是否值得”的讨论
提供了一个可复核的基线。

\appendix
\section{复现}
\label{app:repro}

\subsection{环境与版本}

\begin{itemize}
  \item 操作系统：Linux（开发与门禁面）；桌面端目标为 Windows/Linux/macOS。
  \item 运行时：Flutter/Dart 桌面外壳、Go \texttt{collector}、Rust \texttt{kernel}；
        Python 分析环境为仓库内虚拟环境 \verb|~/.cache/cbr-venv|。
  \item 排版：\texttt{tectonic}（XeTeX 引擎，\texttt{ctex} + Fandol 字体）。
  \item 生成端点：\texttt{\RESULT{meta}{base_url}}；模型
        \texttt{\RESULT{meta}{model}}；\texttt{lock\_models=true}。
  \item 随机种子：\texttt{\RESULT{meta}{seed}}；抽样对数：压力/平静各 60。
  \item 代码提交：\texttt{<COMMIT\_SHA>}。正式归档时用
        \texttt{git rev-parse HEAD} 的结果替换；本稿不手写 SHA，避免与产物不一致。
\end{itemize}

\subsection{命令}

在仓库根目录执行：
\begin{verbatim}
# 1) 构建内核（首次较慢）
cargo build --manifest-path services/kernel/Cargo.toml

# 2) 跑 120 个配对时点，4 分片并行；产物写入 artifacts/bench/shard*.json
PAIRS=60 SHARDS=4 bash tools/bench/run_console_shards.sh

# 3) 合并分片（拒绝重复时点与不一致口径）
~/.cache/cbr-venv/bin/python tools/bench/analysis/merge_shards.py \
    artifacts/bench --out artifacts/bench/console.json

# 4) 评分（输出 precision/recall/lift/Wilson/McNemar）
~/.cache/cbr-venv/bin/python tools/bench/analysis/score_arms.py \
    artifacts/bench/console.json

# 5) 生成论文表格与正文宏
~/.cache/cbr-venv/bin/python tools/bench/analysis/render_tables.py \
    --bench artifacts/bench --out docs/25-paper/tables

# 6) 编译论文（首次会下载 TeX 包；需要代理时设置 HTTPS_PROXY）
HTTPS_PROXY=http://127.0.0.1:7897 \
    /home/xrak/rt/bin/tectonic docs/25-paper/paper.tex
\end{verbatim}

\subsection{产物与字段}

\texttt{artifacts/bench/console.json} 的关键字段：
\begin{itemize}
  \item \texttt{model}、\texttt{base\_url}、\texttt{arms}、\texttt{sampling.seed}：
        实验口径；合并脚本要求所有分片一致。
  \item \texttt{sampling.dates[]}：配对抽样结果，含 \texttt{as\_of}、
        \texttt{pressure} 与 \texttt{y20}。
  \item \texttt{records[].arms.<arm>}：\texttt{risk}、\texttt{ms}、\texttt{tokens}、
        \texttt{partial}、\texttt{claims}、\texttt{conflicts}、\texttt{degraded}、
        \texttt{failure}、\texttt{verification}。
  \item \texttt{verification}：\texttt{passed}、\texttt{numeric\_error\_rate}、
        \texttt{leakage\_events}、\texttt{unsupported\_claims}。
\end{itemize}
\texttt{render\_tables.py} 只读这些字段与黄金向量，不修改产物；缺字段写
\texttt{--}。所有表格文件头部记录数据源与生成时间，正文宏文件
\texttt{tables/macros.tex} 由同一脚本覆盖写入。

\subsection{口径核验与已知差异}

\paragraph{标签口径。} 项目内部报告（\texttt{docs/13} §5.5）曾称 $y_{20}$ 含当日；
经逐日核对，黄金向量中 $y_{20}$ 与 $\max(hs_{t+1},\dots,hs_{t+20})$ 在
\RESULT{label}{n_agree}/\RESULT{label}{n_future} 个可定义日完全一致，即该口径
不含当日，对应条件概率为 \RESULT{label}{p_y_hs1}。该表述已修正，标签生成脚本
同步修正（原脚本用 \texttt{hs.iloc[i:i+20]}，与产物中的严格未来口径不一致）；表~
\ref{tab:circularity} 另给出从 $t+2$ 起算的更严格标签，用于量化自相关折扣。

\paragraph{上一轮数字。} 第 6.1 节引用的 $n=12$、lift 1.09、约 20 秒与 6k token
来自 \texttt{docs/13} 的已发布结果，仅用于解释本轮设计差异；它们不是本轮产物，
也不进入本轮主表。

\paragraph{提交与产物对应。} 本稿的 PDF 由当前工作树编译；正式归档时应把
\texttt{<COMMIT\_SHA>} 替换为产生 \texttt{artifacts/bench} 的提交，并在
\texttt{CHANGELOG} 记录模型端点、提示词版本与抽样种子。若工作树有未提交修改，
应在归档说明中标注 \texttt{dirty}，不得把不同提交的产物混在同一张表里。

\subsection{复现检查清单}

\begin{enumerate}
  \item \texttt{console.json} 存在，且 \texttt{n\_runs} 等于计划时点数；
  \item 三臂的 \texttt{model}/\texttt{base\_url} 一致，\texttt{sampling.seed} 一致；
  \item \texttt{render\_tables.py} 输出中没有把缺失值写成 0；
  \item 主表 $n$ 与 McNemar $n$ 的差异有解释（缺失定级、真值缺失或分片未完成）；
  \item 覆盖率、遮蔽、安慰剂与循环性表均已生成，而不是只生成主表；
  \item PDF 中所有结果数字都能在 \texttt{tables/macros.tex} 中找到对应宏。
\end{enumerate}

\clearpage
\bibliographystyle{unsrtnat}
\bibliography{refs}

\end{document}
